
\section{\BookRecordIndex{T0388}有限维约化}\label{sec:A-4}

本节说明如何为两个巴拿赫空间之间光滑映射 $f:X\to Y$ 的零点集构造局部的
有限维模型。这种\BookRecordIndex{T0388}有限维约化有时称为 \emph{\BookRecordIndex{T0029}Kuranishi 模型}。设 $f(0)=0$，
记 $D:=df(0):X\to Y$；按假设，$D$ 是有界线性算子。$D$ 的一个
\emph{\BookRecordIndex{T0252}拟逆}是有界线性算子 $T:Y\to X$，满足
\[
  TDT=T,\qquad DTD=D.
\]
下一个命题给出\BookRecordIndex{T0252}拟逆存在的充要条件。特别地，每个\BookRecordIndex{T0128}弗雷德霍姆算子都有\BookRecordIndex{T0252}拟逆。

\begin{proposition}\label{prop:A-4-1}
有界线性算子 $D:X\to Y$ 有\BookRecordIndex{T0252}拟逆，当且仅当它满足下列条件：
\begin{enumerate}[label=(\roman*)]
\item $D$ 的像是闭的；
\item $D$ 的核在 $X$ 中有补空间；
\item $D$ 的像在 $Y$ 中有补空间。
\end{enumerate}
\end{proposition}

\begin{proof}
先设 $D$ 满足上述三个条件，记
\[
  X_0:=\ker D,\qquad Y_1:=\operatorname{im}D,
\]
并选择补空间 $X_1,Y_0$，使
\[
  X=X_0\oplus X_1,\qquad Y=Y_0\oplus Y_1.
\]
于是 $X_1,Y_1$ 都是巴拿赫空间，$D$ 在 $X_1$ 上的限制给出双射的有界线性算子
$D_1:X_1\to Y_1$。读者容易验证，对 $y_0\in Y_0$、$y_1\in Y_1$ 定义的
算子 $T:Y\to X$，
\[
  T(y_0+y_1):=D_1^{-1}y_1,
\]
就是 $D$ 的\BookRecordIndex{T0252}拟逆。反过来，若 $T:Y\to X$ 是 $D$ 的\BookRecordIndex{T0252}拟逆，则所需补空间为
\[
  X_1:=\operatorname{im}T,\qquad Y_0:=\ker T.
\]
这是因为
\[
  P:=TD:X\to X,\qquad Q:=DT:Y\to Y
\]
都是投影算子，并且
\[
  \operatorname{im}P=\operatorname{im}T,\quad
  \ker P=\ker D,\quad
  \operatorname{im}Q=\operatorname{im}D,\quad
  \ker Q=\ker T.
\]
这也表明 $D$ 的像是闭的。
\end{proof}

\begin{remark}\label{rem:A-4-2}
设紧李群 $G$ 通过强连续映射作用在巴拿赫空间 $X,Y$ 上，即
\[
  G\to\mathcal L(X):g\mapsto\Phi_g,\qquad
  G\to\mathcal L(Y):g\mapsto\Psi_g.
\]
设 $D:X\to Y$ 是 $G$-等变的有界线性算子。若 $D$ 有\BookRecordIndex{T0252}拟逆，则它也有一个
$G$-等变的\BookRecordIndex{T0252}拟逆。事实上，若 $T$ 是任意\BookRecordIndex{T0252}拟逆，则
\[
  T_g=\Phi_gT\Psi_g^{-1}
\]
也是\BookRecordIndex{T0252}拟逆。因此，取 $G$ 上总体积为 $1$ 的 Haar 测度 $d\mu$，平均算子
\[
  S:=\int_G T_g\,d\mu(g)
\]
是 $G$-等变的，且满足 $DSD=D$。由此，
\[
  R:=SDS=\int_G\int_G T_gDT_h\,d\mu(g)\,d\mu(h)
\]
是 $D$ 的一个 $G$-等变\BookRecordIndex{T0252}拟逆。
\end{remark}

\begin{theorem}\label{thm:A-4-3}
设 $X,Y$ 是巴拿赫空间，$U\subset X$ 是零点的一个邻域，$\ell$ 是正整数。
设 $f:U\to Y$ 是 $C^\ell$ 映射，满足 $f(0)=0$，并设
$D:=df(0):X\to Y$ 有\BookRecordIndex{T0252}拟逆 $T:Y\to X$。则存在零点的开邻域 $W\subset X$、
局部 $C^\ell$ 微分同胚 $g:W\to g(W)$（其中 $g(W)\subset U$ 为开集），
以及 $C^\ell$ 映射 $f_0:W\to Y_0:=\ker T$，使得对 $x\in W$ 有
\[
  f\circ g(x)=f_0(x)+Dx,
\]
并且
\[
  g(0)=0,\qquad dg(0)=\operatorname{id},\qquad
  f_0(0)=0,\qquad df_0(0)=0.
\]
若 $f$ 对某紧李群 $G$ 在 $X,Y$ 上的强连续作用等变，则可以选取
$G$-等变的 $g$ 与 $f_0$。
\end{theorem}

\begin{proof}
考虑 $C^\ell$ 映射 $\psi:U\to X$，定义为
\[
  \psi(x):=x+T(f(x)-Dx).
\]
（这个公式与命题~\ref{prop:A-3-4} 的证明所用公式几乎相同，
只是略去了 $TDx_1$ 一项。）此映射满足
\[
  \psi(0)=0,\qquad d\psi(0)=\operatorname{id}.
\]
由定理~\ref{thm:A-3-1}，$\psi$ 在零点附近有局部逆。
取零点的开邻域 $U_0\subset U$，使 $\psi$ 在其上有逆；记
$W:=\psi(U_0)$，并定义 $g:W\to U_0$ 及
$f_0:W\to\ker T$ 为
\[
  g:=\psi^{-1},\qquad
  f_0:=(\operatorname{id}-DT)\circ f\circ\psi^{-1}.
\]
公式
\[
  D\psi(x)=Dx+DT(f(x)-Dx)=DTf(x)
\]
表明 $D=DT\circ f\circ\psi^{-1}$（等式两侧均视为 $W$ 上的映射），因此
\[
  f\circ g=f\circ\psi^{-1}
  =(\operatorname{id}-DT)\circ f\circ\psi^{-1}+D=f_0+D.
\]
若 $f$ 为 $G$-等变映射，则由注~\ref{rem:A-4-2} 可选取
$G$-等变的\BookRecordIndex{T0252}拟逆 $T$。此时上述公式表明 $g$ 和 $f_0$ 也等变。
\end{proof}

设 $g,f_0$ 如定理~\ref{thm:A-4-3}。则通过 $g$，零点附近 $f$ 的局部
零点集可与 $f_0:W\cap\ker D\to Y_0=\ker T$ 的零点集等同：
\[
  f^{-1}(0)\cap g(W)
  =\bigl\{g(x)\mid x\in W\cap\ker D,\ f_0(x)=0\bigr\}.
\]
要看出这一点，只须注意 $f_0(W)\subset Y_0$，而 $Y_0$ 是
$\operatorname{im}D$ 的补空间，因此
\[
  f_0(x)+Dx=0\quad\Longleftrightarrow\quad Dx=0,\ f_0(x)=0.
\]
当 $D$ 是\BookRecordIndex{T0128}弗雷德霍姆算子时，这一观察尤其有用。此时 $D$ 的核与余核都是有限维的，
因而 $f_0:W\cap\ker D\to Y_0\cong\operatorname{coker}D$ 是有限维向量空间
之间的 $C^\ell$ 映射。
\sourcefootnote{原书末段把 $f_0$ 的定义域印作 $U_0\cap\ker D$，
但定理及其证明定义 $f_0$ 于 $W=\psi(U_0)$；同段前一式也使用 $W\cap\ker D$，
故据此补正。原书此处又称其为 smooth 映射，在仅假设 $f$ 为 $C^\ell$ 时应为
$C^\ell$。}
